\documentclass{beamer}

\usepackage[T1]{fontenc}
\usepackage[utf8]{inputenc}
\usepackage{lmodern}
\usepackage{amsmath,amssymb,amsfonts}
\usepackage{physics}
\usepackage{bm}
\usepackage{mathtools}
\usepackage{quantikz}

\title{The Quantum Approximate Optimization Algorithm (QAOA)}
\subtitle{Theory, Circuits, and Applications}
\author{Morten Hjorth-Jensen}
\date{Spring 2026}

\begin{document}

%------------------------------------------------
\frame{\titlepage}

%------------------------------------------------
\begin{frame}{Outline}
\begin{enumerate}
\item Motivation: Optimization in physics
\item Variational quantum algorithms
\item Definition of QAOA
\item Cost Hamiltonians and problem encoding
\item QAOA circuit construction
\item Example: MaxCut
\item Classical optimization loop
\item Performance and limitations
\end{enumerate}
\end{frame}

%================================================
\begin{frame}{Motivation}
Many problems in physics and computer science reduce to optimization:

\begin{itemize}
\item Ground-state energy minimization
\item Spin models (Ising, Heisenberg)
\item Combinatorial optimization (MaxCut, SAT)
\end{itemize}

General form:
\[
\min_{z \in \{0,1\}^n} C(z)
\]

QAOA is designed to solve such problems using quantum circuits.
\end{frame}

%================================================
\begin{frame}{Variational Quantum Algorithms}
QAOA belongs to the class of hybrid quantum-classical algorithms.

Idea:
\begin{itemize}
\item Prepare a parametrized quantum state \(\ket{\psi(\boldsymbol{\theta})}\)
\item Measure expectation value:
\[
E(\boldsymbol{\theta}) = \bra{\psi(\boldsymbol{\theta})} H_C \ket{\psi(\boldsymbol{\theta})}
\]
\item Optimize parameters classically
\end{itemize}

This is closely related to VQE.
\end{frame}

%================================================
\begin{frame}{Definition of QAOA}
QAOA constructs the state
\[
\ket{\psi_p(\boldsymbol{\gamma}, \boldsymbol{\beta})}
=
\prod_{k=1}^p
e^{-i \beta_k H_M}
e^{-i \gamma_k H_C}
\ket{+}^{\otimes n}
\]

where:
\begin{itemize}
\item \(H_C\): cost Hamiltonian
\item \(H_M\): mixing Hamiltonian
\item \(p\): circuit depth
\end{itemize}
\end{frame}

%================================================
\begin{frame}{Initial State}
The algorithm starts from
\[
\ket{+}^{\otimes n}
=
\frac{1}{2^{n/2}} \sum_z \ket{z}
\]

This represents an equal superposition of all candidate solutions.
\end{frame}

%================================================
\begin{frame}{Cost Hamiltonian}
The classical cost function is encoded as a diagonal Hamiltonian:
\[
H_C = \sum_z C(z) \ket{z}\bra{z}
\]

Example (Ising form):
\[
H_C = \sum_{i<j} J_{ij} Z_i Z_j + \sum_i h_i Z_i
\]

Thus:
\[
C(z) = \bra{z} H_C \ket{z}
\]
\end{frame}

%================================================
\begin{frame}{Mixing Hamiltonian}
The standard choice is
\[
H_M = \sum_{i=1}^n X_i
\]

This generates transitions between computational basis states and ensures exploration of the search space.
\end{frame}

%================================================
\begin{frame}{QAOA Circuit Structure}
Each layer consists of:
\begin{enumerate}
\item Cost unitary:
\[
U_C(\gamma) = e^{-i \gamma H_C}
\]
\item Mixer:
\[
U_M(\beta) = e^{-i \beta H_M}
\]
\end{enumerate}

Repeated \(p\) times.
\end{frame}

%================================================
\begin{frame}{QAOA Circuit (p=1)}
\centering
\begin{quantikz}
\lstick{$\ket{+}$} & \gate{e^{-i\gamma H_C}} & \gate{e^{-i\beta X}} & \meter{} \\
\lstick{$\ket{+}$} & \gate{e^{-i\gamma H_C}} & \gate{e^{-i\beta X}} & \meter{}
\end{quantikz}
\end{frame}

%================================================
\begin{frame}{Example: MaxCut Problem}
Given a graph \(G=(V,E)\), maximize:
\[
C(z) = \sum_{(i,j)\in E} \frac{1 - z_i z_j}{2}
\]

Mapping to Hamiltonian:
\[
H_C = \sum_{(i,j)\in E} \frac{1 - Z_i Z_j}{2}
\]
\end{frame}

%================================================
\begin{frame}{MaxCut Cost Unitary}
For one edge:
\[
e^{-i\gamma (1 - Z_i Z_j)/2}
\]

Circuit implementation:
\centering
\begin{quantikz}
\lstick{$q_i$} & \ctrl{1} & \gate{R_z(2\gamma)} & \ctrl{1} & \qw \\
\lstick{$q_j$} & \targ{}  & \qw                  & \targ{}  & \qw
\end{quantikz}
\end{frame}

%================================================
\begin{frame}{Mixer Unitary}
For each qubit:
\[
e^{-i\beta X}
\]

Circuit:
\[
R_x(2\beta)
\]

\centering
\begin{quantikz}
\lstick{$q_i$} & \gate{R_x(2\beta)} & \qw
\end{quantikz}
\end{frame}

%================================================
\begin{frame}{Expectation Value}
We evaluate:
\[
\langle H_C \rangle =
\bra{\psi(\gamma,\beta)} H_C \ket{\psi(\gamma,\beta)}
\]

This is estimated via measurements on the quantum device.
\end{frame}

%================================================
\begin{frame}{Classical Optimization Loop}
\begin{enumerate}
\item Initialize parameters \((\gamma,\beta)\)
\item Run quantum circuit
\item Measure expectation value
\item Update parameters using classical optimizer
\end{enumerate}

Repeat until convergence.
\end{frame}

%================================================
\begin{frame}{Physical Interpretation}
QAOA can be viewed as a discretized quantum annealing process:

\[
H(t) = (1-s(t))H_M + s(t)H_C
\]

with alternating evolution under \(H_C\) and \(H_M\).
\end{frame}

%================================================
\begin{frame}{Performance}
\begin{itemize}
\item \(p=1\): shallow circuits, limited accuracy
\item larger \(p\): better approximation
\item classical optimization becomes harder
\end{itemize}

Tradeoff:
\[
\text{accuracy} \leftrightarrow \text{circuit depth}
\]
\end{frame}

%================================================
\begin{frame}{Relation to VQE}
QAOA is a special case of VQE:
\begin{itemize}
\item structured ansatz
\item problem-inspired Hamiltonians
\end{itemize}

Difference:
\begin{itemize}
\item QAOA: alternating operators
\item VQE: general parametrized circuits
\end{itemize}
\end{frame}

%================================================
\begin{frame}{Limitations}
\begin{itemize}
\item Barren plateaus in optimization
\item Hardware noise
\item scaling with problem size
\end{itemize}
\end{frame}

%================================================
\begin{frame}{Summary}
QAOA:
\begin{itemize}
\item encodes optimization problems into Hamiltonians
\item uses alternating quantum evolution
\item is hybrid quantum-classical
\item connects quantum computing with statistical physics
\end{itemize}
\end{frame}



%================================================
% PART II: Adiabatic Origins and Analytics
%================================================

%------------------------------------------------
\frame{\titlepage}

%------------------------------------------------
\begin{frame}{Outline}
\begin{enumerate}
\item QAOA review
\item Adiabatic quantum computation
\item Trotterization $\Rightarrow$ QAOA
\item MaxCut as an Ising model
\item Analytical solution at $p=1$
\item Statistical mechanics interpretation
\end{enumerate}
\end{frame}

%================================================
\begin{frame}{QAOA State}
\[
\ket{\psi_p(\boldsymbol{\gamma},\boldsymbol{\beta})}
=
\prod_{k=1}^p e^{-i\beta_k H_M} e^{-i\gamma_k H_C}
\ket{+}^{\otimes n}
\]

\begin{itemize}
\item $H_C$: problem Hamiltonian
\item $H_M = \sum_i X_i$: mixing Hamiltonian
\end{itemize}

Goal:
\[
\max \bra{\psi} H_C \ket{\psi}
\]
\end{frame}

%================================================
\begin{frame}{Adiabatic Quantum Computation}
Consider a time-dependent Hamiltonian:
\[
H(s) = (1-s)H_M + s H_C, \quad s \in [0,1]
\]

Adiabatic theorem:
\begin{itemize}
\item If evolution is slow, system stays in ground state
\item Final state $\rightarrow$ ground state of $H_C$
\end{itemize}

Time evolution:
\[
U = \mathcal{T}\exp\left(-i\int_0^T H(s) ds\right)
\]
\end{frame}

%================================================
\begin{frame}{Discretization of Adiabatic Evolution}
Divide time into $p$ steps:
\[
U \approx \prod_{k=1}^p
e^{-i \Delta t ((1-s_k)H_M + s_k H_C)}
\]

Using Trotter decomposition:
\[
e^{-i (A+B)\Delta t}
\approx
e^{-iA\Delta t} e^{-iB\Delta t}
\]
\end{frame}

%================================================
\begin{frame}{From Adiabatic Evolution to QAOA}
Applying Trotterization:
\[
U \approx \prod_{k=1}^p
e^{-i\beta_k H_M} e^{-i\gamma_k H_C}
\]

Thus:
\[
\boxed{\text{QAOA is a discretized (Trotterized) adiabatic evolution}}
\]

Key idea:
\begin{itemize}
\item Replace fixed schedule with variational parameters
\end{itemize}
\end{frame}

%================================================
\begin{frame}{MaxCut as an Ising Model}
MaxCut cost:
\[
C(z) = \sum_{(i,j)\in E} \frac{1 - z_i z_j}{2}
\]

Hamiltonian form:
\[
H_C = \sum_{(i,j)\in E} \frac{1 - Z_i Z_j}{2}
\]

Equivalent to Ising model:
\[
H = -\sum_{i<j} J_{ij} Z_i Z_j
\]

Thus QAOA solves an Ising ground-state problem.
\end{frame}

%================================================
\begin{frame}{Statistical Mechanics Interpretation}
Define partition function:
\[
Z = \mathrm{Tr}(e^{-\beta H})
\]

QAOA instead prepares:
\[
\ket{\psi(\gamma,\beta)} \sim e^{-i\beta H_M} e^{-i\gamma H_C}
\]

Interpretation:
\begin{itemize}
\item Imaginary time: classical Gibbs states
\item Real time: quantum interference
\end{itemize}

QAOA interpolates between:
\[
\text{thermal sampling} \leftrightarrow \text{coherent dynamics}
\]
\end{frame}

%================================================
\begin{frame}{Analytical Example: Single Edge}
Consider two qubits with one edge:
\[
H_C = \frac{1 - Z_1 Z_2}{2}
\]

Initial state:
\[
\ket{+}^{\otimes 2}
\]
\end{frame}

%================================================
\begin{frame}{Apply Cost Unitary}
\[
U_C(\gamma) = e^{-i\gamma H_C}
\]

Since $Z_1 Z_2$ eigenvalues $\pm1$:
\[
U_C = e^{-i\gamma/2} e^{i\gamma Z_1 Z_2 /2}
\]
\end{frame}

%================================================
\begin{frame}{Apply Mixer}
\[
U_M(\beta) = e^{-i\beta (X_1 + X_2)}
\]

Acts as rotations:
\[
R_x(2\beta)
\]
on each qubit.
\end{frame}

%================================================
\begin{frame}{Expectation Value}
We compute:
\[
\langle H_C \rangle =
\frac{1}{2}(1 - \langle Z_1 Z_2 \rangle)
\]

After algebra (standard result):
\[
\langle Z_1 Z_2 \rangle =
\sin(4\beta)\sin(2\gamma)
\]
\end{frame}

%================================================
\begin{frame}{Optimization (p=1)}
Thus:
\[
\langle H_C \rangle =
\frac{1}{2}(1 - \sin(4\beta)\sin(2\gamma))
\]

Maximize by choosing:
\[
\beta = \frac{\pi}{8}, \quad \gamma = \frac{\pi}{4}
\]

This gives optimal cut value.
\end{frame}

%================================================
\begin{frame}{Interpretation of p=1 Result}
Key insight:
\begin{itemize}
\item Even shallow circuits capture nontrivial correlations
\item Performance improves with connectivity
\end{itemize}

For general graphs:
\begin{itemize}
\item expectation depends on local graph structure
\end{itemize}
\end{frame}

%================================================
\begin{frame}{Connection to Spin Systems}
QAOA operates directly on spin Hamiltonians:
\[
H = \sum J_{ij} Z_i Z_j
\]

This is:
\begin{itemize}
\item classical Ising model
\item quantum transverse-field Ising model
\end{itemize}

Thus QAOA = controlled exploration of spin configurations.
\end{frame}

%================================================
\begin{frame}{QAOA vs Classical Methods}
Classical:
\begin{itemize}
\item simulated annealing
\item Monte Carlo
\end{itemize}

QAOA:
\begin{itemize}
\item uses interference instead of thermal sampling
\item explores many configurations simultaneously
\end{itemize}
\end{frame}

%================================================
\begin{frame}{Summary}
\begin{itemize}
\item QAOA arises from Trotterized adiabatic evolution
\item Encodes optimization as Ising Hamiltonians
\item p=1 case admits analytical solutions
\item Connects quantum computing with statistical physics
\end{itemize}
\end{frame}

%================================================
% PART III: Statistical Physics and Control Theory
%================================================

%------------------------------------------------
\frame{\titlepage}

%------------------------------------------------
\begin{frame}{Overview}
\begin{itemize}
\item Optimization problems and spin systems
\item QAOA construction
\item Adiabatic origin and Trotterization
\item Analytical structure (p=1)
\item General graphs and locality
\item Landscape and optimization theory
\item Statistical mechanics connections
\item Control-theoretic interpretation
\end{itemize}
\end{frame}

%================================================
\begin{frame}{Optimization as Physics}
Many optimization problems map to:
\[
\min_z C(z)
\]

Equivalent to finding ground state of:
\[
H_C = \sum_z C(z)\ket{z}\bra{z}
\]

This is identical to:
\begin{itemize}
\item classical spin systems
\item Ising Hamiltonians
\end{itemize}
\end{frame}

%================================================
\begin{frame}{Ising Representation}
General Ising form:
\[
H_C = \sum_{i<j} J_{ij} Z_i Z_j + \sum_i h_i Z_i
\]

Binary variables:
\[
z_i \in \{-1,1\}
\]

Thus QAOA operates directly on many-body Hamiltonians.
\end{frame}

%================================================
\begin{frame}{QAOA Ansatz}
\[
\ket{\psi_p(\gamma,\beta)} =
\prod_{k=1}^p e^{-i\beta_k H_M} e^{-i\gamma_k H_C}
\ket{+}^{\otimes n}
\]

Mixer:
\[
H_M = \sum_i X_i
\]
\end{frame}

%================================================
\begin{frame}{Physical Picture}
\begin{itemize}
\item $H_C$: encodes energy landscape
\item $H_M$: induces transitions
\end{itemize}

Thus:
\[
\text{QAOA = controlled quantum exploration of configuration space}
\]
\end{frame}

%================================================
\begin{frame}{Circuit Representation}
\centering
\begin{quantikz}
\lstick{$\ket{+}$} & \gate{e^{-i\gamma H_C}} & \gate{e^{-i\beta X}} & \qw \\
\lstick{$\ket{+}$} & \gate{e^{-i\gamma H_C}} & \gate{e^{-i\beta X}} & \qw
\end{quantikz}
\end{frame}

%================================================
\begin{frame}{Adiabatic Evolution}
\[
H(s) = (1-s)H_M + s H_C
\]

Evolution:
\[
U = \mathcal{T} e^{-i\int_0^T H(s)ds}
\]
\end{frame}

%================================================
\begin{frame}{Trotterization}
Discretize:
\[
U \approx \prod_{k=1}^p e^{-i\Delta t H(s_k)}
\]

Apply Trotter:
\[
e^{-i(A+B)} \approx e^{-iA} e^{-iB}
\]
\end{frame}

%================================================
\begin{frame}{QAOA as Trotterized Adiabatic Evolution}
\[
\boxed{
U \approx \prod_{k=1}^p e^{-i\beta_k H_M} e^{-i\gamma_k H_C}
}
\]

Key difference:
\begin{itemize}
\item adiabatic: fixed schedule
\item QAOA: variational parameters
\end{itemize}
\end{frame}

%================================================
\begin{frame}{MaxCut Problem}
\[
C(z)=\sum_{(i,j)} \frac{1-z_i z_j}{2}
\]

Hamiltonian:
\[
H_C = \sum_{(i,j)} \frac{1-Z_i Z_j}{2}
\]
\end{frame}

%================================================
\begin{frame}{Local Structure}
Important fact:
\begin{itemize}
\item At depth $p$, correlations extend only distance $p$
\end{itemize}

Thus expectation values depend only on local subgraphs.
\end{frame}

%================================================
\begin{frame}{Analytical p=1: Setup}
Single edge:
\[
H_C = \frac{1-Z_1Z_2}{2}
\]

Initial state:
\[
\ket{+}^{\otimes 2}
\]
\end{frame}

%================================================
\begin{frame}{Evolution}
State:
\[
\ket{\psi} = e^{-i\beta H_M} e^{-i\gamma H_C} \ket{++}
\]
\end{frame}

%================================================
\begin{frame}{Key Computation}
One finds:
\[
\langle Z_1 Z_2 \rangle = \sin(4\beta)\sin(2\gamma)
\]

Thus:
\[
\langle H_C \rangle = \frac{1}{2}(1-\sin(4\beta)\sin(2\gamma))
\]
\end{frame}

%================================================
\begin{frame}{Optimal Parameters}
Maximized at:
\[
\beta = \frac{\pi}{8}, \quad \gamma = \frac{\pi}{4}
\]
\end{frame}

%================================================
\begin{frame}{General Graphs}
For arbitrary graph:
\begin{itemize}
\item expectation depends on local neighborhoods
\item can be computed analytically for small $p$
\end{itemize}
\end{frame}

%================================================
\begin{frame}{Statistical Mechanics View}
Define:
\[
Z = \mathrm{Tr}(e^{-\beta H})
\]

QAOA:
\[
\ket{\psi} = e^{-i\beta H_M} e^{-i\gamma H_C}\ket{+}
\]
\end{frame}

%================================================
\begin{frame}{Real vs Imaginary Time}
\begin{itemize}
\item imaginary time: $e^{-\beta H}$ → Gibbs states
\item real time: $e^{-iHt}$ → unitary evolution
\end{itemize}

QAOA operates in real time.
\end{frame}

%================================================
\begin{frame}{Relation to Annealing}
\begin{itemize}
\item simulated annealing: thermal fluctuations
\item quantum annealing: tunneling
\item QAOA: controlled coherent evolution
\end{itemize}
\end{frame}

%================================================
\begin{frame}{Energy Landscape}
Define:
\[
E(\gamma,\beta)=\bra{\psi}H_C\ket{\psi}
\]

Properties:
\begin{itemize}
\item highly non-convex
\item periodic
\end{itemize}
\end{frame}

%================================================
\begin{frame}{Parameter Landscape}
\begin{itemize}
\item small $p$: smooth landscape
\item large $p$: complex landscape
\end{itemize}

Connections to:
\begin{itemize}
\item spin-glass landscapes
\end{itemize}
\end{frame}

%================================================
\begin{frame}{Control-Theoretic View}
QAOA:
\[
U(\theta)=\mathcal{T} e^{-i\int dt [u_1(t)H_C + u_2(t)H_M]}
\]

Piecewise constant controls:
\[
u_1(t), u_2(t)
\]
\end{frame}

%================================================
\begin{frame}{Continuous Limit}
As $p\to\infty$:
\[
\text{QAOA} \to \text{continuous Schr\"odinger evolution}
\]
\end{frame}

%================================================
\begin{frame}{Connection to Optimal Control}
Goal:
\[
\max \bra{\psi(T)}H_C\ket{\psi(T)}
\]

This is a quantum control problem.
\end{frame}

%================================================
\begin{frame}{Relation to VQE}
\begin{itemize}
\item QAOA: structured ansatz
\item VQE: general ansatz
\end{itemize}
\end{frame}

%================================================
\begin{frame}{Scaling Challenges}
\begin{itemize}
\item parameter optimization
\item noise
\item circuit depth
\end{itemize}
\end{frame}

%================================================
\begin{frame}{Barren Plateaus}
Gradients may vanish exponentially:
\[
\partial_\theta E \to 0
\]
\end{frame}

%================================================
\begin{frame}{Why QAOA Still Works}
Structure:
\begin{itemize}
\item locality
\item problem-inspired ansatz
\end{itemize}
\end{frame}

%================================================
\begin{frame}{Summary}
\begin{itemize}
\item QAOA = Trotterized adiabatic evolution
\item maps optimization → spin systems
\item admits analytic results at small $p$
\item connects quantum computing and statistical physics
\end{itemize}
\end{frame}

\end{document}
